% Fragmento de trabajo: propuesta de introducción para el plan de tesis.
% Mantiene la cita autor-año usada por el proyecto. Es un borrador independiente;
% no reemplaza automáticamente partes/04AspectosGenerales.tex.

\subsection{Introducción}

El procesamiento del lenguaje natural (PLN) permite analizar y generar textos mediante modelos computacionales, pero la utilidad de estas tecnologías depende también de la disponibilidad de datos y recursos lingüísticos adecuados. Esta condición es relevante para el Quechua Collao, cuya presencia en herramientas digitales requiere corpus y recursos que representen la variedad. Entre los antecedentes incluidos en este plan figura la construcción de un corpus paralelo Collao quechua--español a partir de textos de la Academia Mayor de la Lengua Quechua; también se han desarrollado recursos léxicos y morfológicos para quechua \citep{Orcotoma2025,Duran2026}. Estos aportes fortalecen la base de recursos, aunque no equivalen por sí solos a un sistema de generación contextualizada que funcione localmente.

Entre los enfoques posibles se encuentran los modelos de lenguaje que reciben contexto lingüístico recuperado, los modelos compactos ejecutados en dispositivos locales y los sistemas de generación aumentada por recuperación (RAG), que incorporan documentos externos durante la respuesta. En traducción de Quechua sureño al español, Court y Elsner \citeyearpar{Court2024} muestran que el uso del contexto depende del tipo y la pertinencia de la información proporcionada. En otro dominio, una evaluación de RAG offline con modelos pequeños reporta dificultades cuando el contexto es extenso o contiene fragmentos irrelevantes \citep{Hevia2025}. Estos resultados corresponden a tareas y datos diferentes de la generación en Quechua Collao; por tanto, motivan una evaluación específica en lugar de asumir que el RAG siempre mejora o siempre perjudica la respuesta. Para ejecutar el sistema sin conexión se plantea un modelo compacto de la familia Qwen2.5 y una base documental local \citep{qwen2024,Xu2024}. La adaptación eficiente mediante QLoRA se considera una alternativa de diseño cuya contribución debe medirse experimentalmente \citep{Dettmers2023}.

Los antecedentes que actualmente recoge el plan cubren la creación de recursos para Quechua Collao, la consulta de información lingüística para tareas de traducción y el uso de modelos pequeños con RAG offline en otros contextos. Sin embargo, esos trabajos no constituyen por sí mismos una evaluación conjunta de un modelo compacto, recuperación documental local y adaptación QLoRA para generar respuestas contextualizadas en Quechua Collao bajo restricciones de ejecución offline. En consecuencia, el vacío se delimita provisionalmente a determinar el comportamiento de esa combinación. Antes de presentar una afirmación concluyente de novedad, esta delimitación debe contrastarse con una búsqueda bibliográfica actualizada y con las condiciones, métricas y limitaciones de los estudios relacionados.

A partir de este vacío, la investigación examinará cómo influye el diseño e implementación de una arquitectura SLM--RAG con QLoRA en la generación contextualizada en Quechua Collao. La comparación propuesta considera el modelo base, el sistema con RAG estándar y la configuración integrada con adaptación, bajo condiciones de inferencia offline documentadas. El interés está en observar la fidelidad lingüística, el anclaje de las respuestas en los documentos recuperados y el costo computacional, sin anticipar de antemano cuál configuración obtendrá mejores resultados. La pregunta general correspondiente se formula en el apartado de formulación del problema.

La propuesta consiste en diseñar e implementar la arquitectura integrada con un modelo compacto cuantizado, recuperación local de documentos y adaptadores QLoRA. El objetivo general conserva la formulación del plan: diseñar e implementar una arquitectura SLM--RAG con QLoRA para la generación contextualizada en Quechua Collao dentro de un entorno computacional offline. Su valor científico será aportar evidencia sobre la interacción entre recuperación y adaptación para esta variedad; su valor metodológico, establecer una comparación reproducible entre configuraciones; y su valor tecnológico, determinar las condiciones y límites de funcionamiento local. El objetivo y la justificación se desarrollan en sus apartados respectivos para evitar repetirlos extensamente aquí.

La evaluación se plantea como una comparación experimental entre el SLM base, el SLM con RAG estándar y la arquitectura con RAG y QLoRA. Se informarán las métricas de calidad lingüística y factualidad definidas en el plan, junto con latencia y uso de memoria. La introducción anticipa este enfoque; los conjuntos de datos, parámetros, equipos, procedimientos y análisis se detallan en la sección metodológica.

La memoria final se organiza tentativamente en seis capítulos, según la estructura consignada en el plan: el capítulo I presenta la introducción; el capítulo II reúne el marco teórico y el estado del arte; el capítulo III describe la metodología y el diseño de la arquitectura; el capítulo IV presenta la implementación computacional; el capítulo V expone los resultados y la evaluación experimental; y el capítulo VI desarrolla las conclusiones y los trabajos futuros.
